\documentclass[11pt,a4paper]{article}

\usepackage[T1]{fontenc}
\usepackage[utf8]{inputenc}
\usepackage{lmodern}
\usepackage[margin=2.35cm]{geometry}
\usepackage{amsmath,amssymb,mathtools}
\usepackage{booktabs,longtable,array}
\usepackage[dvipsnames]{xcolor}
\usepackage{hyperref}
\usepackage{microtype}

\hypersetup{
  colorlinks=true,
  hypertexnames=false,
  linkcolor=MidnightBlue,
  urlcolor=MidnightBlue,
  citecolor=MidnightBlue,
  pdftitle={Complete Mathematical Specification of the MO-OBAM PRTI-GH Work}
}

\newcommand{\R}{\mathbb{R}}
\newcommand{\N}{\mathbb{N}}
\newcommand{\I}{\mathbb{I}}
\newcommand{\card}[1]{\left\lvert #1\right\rvert}
\newcommand{\mean}[1]{\overline{#1}}
\newcommand{\given}{\,\middle\vert\,}
\newcommand{\eps}{\varepsilon}
\newcommand{\code}[1]{\texttt{#1}}
\newcolumntype{L}[1]{>{\raggedright\arraybackslash}p{#1}}
\DeclareMathOperator*{\argmin}{arg\,min}
\DeclareMathOperator*{\argmax}{arg\,max}
\DeclareMathOperator{\median}{median}
\DeclareMathOperator{\mode}{mode}
\DeclareMathOperator{\Var}{Var}
\DeclareMathOperator{\KL}{KL}
\DeclareMathOperator{\JSD}{JSD}

\title{Complete Mathematical Specification\\
of the MO-OBAM PRTI-GH Work}
\author{Repository-grounded mathematical compendium}
\date{11 November 2026}

\begin{document}
\maketitle

\begin{abstract}
This document collects the scientific and analytical mathematics actually used
by the repository: data representation, the MO-OBAM objective and constraints,
the exact and particle-swarm solvers, the Pareto extension, disclosure-risk and
diversity measures, before--after distribution statistics, cross-seed summaries,
screening gates, configuration selection, held-out selection, and experiment-size
arithmetic.  It distinguishes the paper equations from project-defined evaluation
and selection rules.  Cryptographic hash internals, file-locking, spreadsheet
formatting, and plotting coordinates are engineering mechanisms and are outside
the scientific-mathematics scope.
\end{abstract}

\tableofcontents
\newpage

\section{Status labels and scope}

Four labels are used throughout:
\begin{itemize}
  \item \textbf{MO-OBAM equation}: a mathematical definition implemented from
  the active arXiv v2 formulation.
  \item \textbf{Optimizer mechanic}: mathematics used by the scalar PSO or the
  additive Pareto PSO extension.
  \item \textbf{Project metric}: a post-anonymization measurement implemented
  for PRTI-GH evaluation.
  \item \textbf{Project rule}: an eligibility, ranking, or held-out rule defined
  by this project rather than asserted as a universal privacy standard.
\end{itemize}

Every displayed formula has one of the following implementation provenances:
\begin{itemize}
  \item \textbf{Direct}: the formula is the literal scalar or array arithmetic in
  the repository.
  \item \textbf{Equivalent}: the formula is an algebraically exact statement of
  a loop, grouping, vectorized operation, comparison, or stable sort in the code.
  \item \textbf{Delegated}: the repository calls NumPy, pandas, or SciPy, and the
  displayed formula is the definition implemented by that invoked library routine.
  \item \textbf{Configured arithmetic}: the formula formalizes a declared grid,
  threshold, gate, or count that the repository constructs programmatically.
\end{itemize}
Thus ``coming from the code'' does not mean that every displayed formula occurs as
one Python source line.  It means that every formula is executed directly, is
algebraically identical to executed operations, defines an invoked library result,
or formalizes an executed project protocol.  No alternative scientific formula is
introduced solely for exposition.

The primary PRTI-GH model has
\[
n=1485,\qquad p_N=2,\qquad p_C=27,\qquad p_Q=p_N+p_C=29,
\qquad p_S=77.
\]
Here $p_N$, $p_C$, $p_Q$, and $p_S$ denote the numbers of numerical QIDs,
categorical QIDs, all QIDs, and sensitive attributes (SAs), respectively.

\section{Notation and data representation}

\subsection{Index sets and variables}

Let
\[
\mathcal I=\{1,\ldots,n\},\quad
\mathcal C=\{1,\ldots,C\},\quad
\mathcal J_N=\{1,\ldots,p_N\},
\]
\[
\mathcal J_C=\{p_N+1,\ldots,p_N+p_C\},\quad
\mathcal J_S=\{1,\ldots,p_S\}.
\]
For record $i$, $x_{ij}$ is QID $j$, $S_{ij}$ is SA $j$, $q_{cj}$ is the
representative of QID $j$ in cluster $c$, and
\[
w_{ic}\in\{0,1\}
\]
is the record--cluster membership variable.  The model hyperparameters are
minimum cluster size $k$, configured cluster count $C$, and scalar entropy
weight $\lambda\in[0,1]$.

\subsection{Deterministic categorical coding}

For each categorical attribute $j$, the preparation layer constructs a
one-to-one code map
\[
e_j:\mathcal V_j\longrightarrow\{1,\ldots,\card{\mathcal V_j}\},
\qquad e_j(v_1)=e_j(v_2)\iff v_1=v_2.
\]
Missing categorical values receive code $0$.  The codes are identifiers, not
ordinal magnitudes.  Consequently, categorical distance uses equality only:
\[
\delta(a,b)=\I(a\ne b)=
\begin{cases}
1,&a\ne b,\\
0,&a=b.
\end{cases}
\tag{1}
\]

For an attribute with $d_j$ distinct non-missing states and $m_j$ missing
records, the audit reports
\[
\operatorname{cardinalityFraction}_j=\frac{d_j}{n},
\qquad
\operatorname{missingFraction}_j=\frac{m_j}{n}.
\tag{2}
\]
For state $v$, its mapping-table proportion is
\[
\widehat p_j(v)=\frac{\#\{i:x_{ij}=v\}}{n}.
\tag{3}
\]

\subsection{Lossless preparation check}

Let $f_j^{\mathrm{src}}(v)$ and $f_j^{\mathrm{dec}}(v)$ be the exact
source and decoded frequencies of state $v$, treating missingness and data
types as part of the state identity.  The largest round-trip discrepancy is
\[
\Delta_j^{\mathrm{freq}}
=\max_v\left|f_j^{\mathrm{src}}(v)-f_j^{\mathrm{dec}}(v)\right|.
\tag{4}
\]
The implemented preparation invariant is
\[
\Delta_j^{\mathrm{freq}}=0\qquad\text{for every attribute }j,
\tag{5}
\]
together with equal shape, column order, and cellwise values after decoding.

\section{MO-OBAM mathematical formulation}

\subsection{Exactly one cluster per record: paper Equation (9)}

\begin{equation}
\sum_{c=1}^{C}w_{ic}=1,
\qquad i\in\mathcal I.
\tag{P9}
\end{equation}

The implementation also enforces $w_{ic}\in\{0,1\}$.  Assignments may
therefore be represented by labels
\[
a_i=\argmax_{c\in\mathcal C}w_{ic},
\qquad w_{ic}=\I(a_i=c).
\tag{6}
\]

\subsection{Numerical and categorical deviations: paper Equations (10)--(13)}

The numerical deviation within cluster $c$ is
\begin{equation}
D^{(N)}_c=
\sum_{i=1}^{n}w_{ic}\sum_{j\in\mathcal J_N}(x_{ij}-q_{cj})^2.
\tag{P10}
\end{equation}
The categorical deviation is
\begin{equation}
D^{(C)}_c=
\sum_{i=1}^{n}w_{ic}\sum_{j\in\mathcal J_C}\delta(x_{ij},q_{cj}).
\tag{P11}
\end{equation}
The record-to-centroid distance is
\begin{equation}
d(i,c)=
\sum_{j\in\mathcal J_N}(x_{ij}-q_{cj})^2
+\sum_{j\in\mathcal J_C}\delta(x_{ij},q_{cj}).
\tag{P12}
\end{equation}
The total deviation introduced by anonymization is
\begin{equation}
D_{\mathrm{anon}}
=\sum_{c=1}^{C}\sum_{i=1}^{n}w_{ic}d(i,c)
=\sum_{c=1}^{C}\left(D^{(N)}_c+D^{(C)}_c\right).
\tag{P13}
\end{equation}

Nearest-centroid assignment in both PSO solvers is
\[
a_i=\argmin_{c\in\mathcal C}d(i,c).
\tag{7}
\]
The array implementation uses the first minimum in a distance tie.

\subsection{Original deviation and information loss: paper Equation (14)}

Let
\[
\overline X_j=\frac1n\sum_{i=1}^{n}x_{ij}
\quad(j\in\mathcal J_N)
\]
and let $\widehat X_j$ be the first observed state among the most frequent
states of categorical QID $j$.  The denominator is
\begin{equation}
D_{\mathrm{original}}=
\sum_{j\in\mathcal J_N}\sum_{i=1}^{n}(x_{ij}-\overline X_j)^2
+\sum_{j\in\mathcal J_C}\sum_{i=1}^{n}
\delta(x_{ij},\widehat X_j).
\tag{P14}
\end{equation}
Information loss is
\begin{equation}
IL=\frac{D_{\mathrm{anon}}}{D_{\mathrm{original}}}.
\tag{IL}
\end{equation}
The implementation requires $D_{\mathrm{original}}>0$ and does not add a
smoothing constant.  Although $IL\ge0$, it is not constrained to be at most
one; a poor set of searched representatives can give $IL>1$.

\subsection{Entropy over multiple sensitive attributes: paper Equations (15)--(19)}

The size of cluster $c$ is
\[
n_c=\sum_{i=1}^{n}w_{ic}.
\tag{$n_c$}
\]
For sensitive attribute $j$ and cluster $c$, the paper entropy is
\begin{equation}
H_{cj}=-\sum_{s:p_{cjs}>0}p_{cjs}\log_2 p_{cjs}.
\tag{P15}
\end{equation}
For state $s$, its within-cluster probability is
\begin{equation}
p_{cjs}
=\frac{\sum_{i=1}^{n}w_{ic}\I(S_{ij}=s)}
{\sum_{i=1}^{n}w_{ic}}.
\tag{P16}
\end{equation}
The indicator is
\[
\I(S_{ij}=s)=
\begin{cases}
1,&S_{ij}=s,\\
0,&\text{otherwise}.
\end{cases}
\tag{P17}
\]
The worst-cluster entropy for SA $j$ and the total entropy term are
\begin{equation}
H_j^{\min}=\min_{c\in\mathcal C}H_{cj},
\tag{P18}
\end{equation}
\begin{equation}
H_{\mathrm{sum}}=\sum_{j=1}^{p_S}H_j^{\min}.
\tag{P19}
\end{equation}
Missing SA values are explicit states and therefore enter the probabilities
and entropy in exactly the same way as other observed states.  The sum is
unweighted: every one of the 77 SAs contributes one minimum-entropy term.

\subsection{Hard \texorpdfstring{$k$}{k}-anonymity: paper Equation (20)}

\begin{equation}
\sum_{i=1}^{n}w_{ic}=n_c\ge k,
\qquad c\in\mathcal C.
\tag{P20}
\end{equation}
This applies to every configured cluster, including one that would otherwise
be empty.  A necessary mathematical feasibility condition is
\begin{equation}
kC\le n,
\qquad
C\le\left\lfloor\frac nk\right\rfloor.
\tag{23}
\end{equation}
Failure of a stochastic initialization to discover a feasible particle is not
the same statement as violation of the necessary condition in (23), nor is it
a proof of mathematical infeasibility.

\subsection{Active scalar objective}

The active scalar problem is
\begin{equation}
\min_{q,w}\;F_{\lambda}(q,w)
=IL(q,w)-\lambda H_{\mathrm{sum}}(w),
\qquad 0\le\lambda\le1,
\tag{24}
\end{equation}
subject to the assignment and hard-$k$ constraints and using the deviations
and entropy above.  A negative objective is valid.  Objective values at
different $\lambda$ values do not share a fixed scale and are not ranked
directly across $\lambda$ in the project selection rules.

\subsection{Anonymization transformation}

Only QIDs are replaced.  For every QID $j$,
\begin{equation}
x'_{ij}=q_{a_i j};
\tag{25}
\end{equation}
for every sensitive or otherwise non-QID attribute,
\begin{equation}
S'_{ij}=S_{ij}.
\tag{26}
\end{equation}

\section{Centroid mathematics and exact small-data validation}

\subsection{Best representatives for fixed membership}

For a fixed nonempty cluster $c$, squared error is minimized by the arithmetic
mean:
\begin{equation}
q^*_{cj}=\frac{1}{n_c}\sum_{i:a_i=c}x_{ij},
\qquad j\in\mathcal J_N.
\tag{27}
\end{equation}
Categorical mismatch is minimized by any within-cluster mode:
\begin{equation}
q^*_{cj}\in\argmax_{v}\sum_{i:a_i=c}\I(x_{ij}=v),
\qquad j\in\mathcal J_C.
\tag{28}
\end{equation}
If modes tie, the exact solver chooses the first tied mode observed in record
order.  Entropy depends on $w$, not $q$, so (27)--(28) minimize the scalar
objective over $q$ for fixed membership.

\subsection{Exhaustive partition search}

For small validation data, the exact solver enumerates every unlabeled
partition of $n$ records into exactly $C$ nonempty blocks satisfying
$n_c\ge k$.  Restricted-growth labels enforce
\begin{equation}
a_1=0,
\qquad
a_i\le1+\max_{h<i}a_h,
\tag{29}
\end{equation}
so relabelings of the same partition are not evaluated repeatedly.  At a
partial assignment, suppose exactly $p$ records have been assigned.  The code's
next zero-based array position is then also $p$, and $r=n-p$ records remain.  The
branch is pruned when
\begin{equation}
\sum_{c=1}^{C}\max(k-n_c,0)>r,
\qquad r=n-p,
\tag{30}
\end{equation}
because the remaining records cannot fill all cluster deficits.  The globally
best enumerated objective is
\begin{equation}
(w^*,q^*)\in\argmin_{(w,q)\in\mathcal F}F_\lambda(w,q),
\tag{31}
\end{equation}
where $\mathcal F$ is the complete feasible partition set and $q$ is given by
(27)--(28).  This exact route is intentionally limited to at most 12 records.

\section{Lossless QID-profile acceleration}

Let $u=1,\ldots,U$ index distinct QID tuples, let $m_u$ be the number of
records with profile $u$, and let $z_u$ be that profile.  Since all records in
a profile have the same distance to every centroid,
\begin{equation}
a_u=\argmin_{c\in\mathcal C}d(z_u,q_c),
\qquad a_i=a_{u(i)}.
\tag{32}
\end{equation}
Cluster sizes and anonymization deviation can therefore be evaluated as
\begin{equation}
n_c=\sum_{u=1}^{U}m_u\I(a_u=c),
\tag{33}
\end{equation}
\begin{equation}
D_{\mathrm{anon}}=\sum_{u=1}^{U}m_u d(z_u,q_{a_u}).
\tag{34}
\end{equation}
For SA $j$ and state $s$, let $m_{ujs}$ be the count of that state within QID
profile $u$.  The cluster/state count used in entropy is
\begin{equation}
N_{cjs}=\sum_{u=1}^{U}m_{ujs}\I(a_u=c),
\qquad p_{cjs}=\frac{N_{cjs}}{n_c}.
\tag{35}
\end{equation}
Equations (32)--(35) are algebraically identical to record-level evaluation.

The compiled categorical kernel returns the Hamming fraction.  With $p_C$
categorical QIDs, exact mismatch count is recovered as
\begin{equation}
d_C(z_u,q_c)=\operatorname{round}\!\left(
p_C\,d_{\mathrm{Hamming}}(z_u,q_c)\right).
\tag{36}
\end{equation}

\section{Scalar particle-swarm optimization}

\subsection{Particle representation and initialization}

Particle $r$ contains all centroid coordinates
\[
Q_r^{(t)}=\{q_{rcj}^{(t)}:c\in\mathcal C,
j\in\mathcal J_N\cup\mathcal J_C\}.
\]
Each initial numerical centroid coordinate is sampled uniformly from the
distinct observed values of its QID, and each categorical coordinate is
sampled uniformly from its encoded category levels.  Initial velocities are
zero.

\subsection{Continuous update}

For each particle and iteration, the implementation draws broadcast
coefficients
\begin{equation}
\alpha_{r,t},\beta_{r,t}\overset{\mathrm{ind}}{\sim}
\operatorname{Uniform}(0,r_{\max}),
\qquad r_{\max}=0.01
\tag{37}
\end{equation}
in the recorded studies, and sets
\begin{equation}
\omega_{r,t}=1-\alpha_{r,t}-\beta_{r,t}>0.
\tag{38}
\end{equation}
For a numerical coordinate, with personal best $P_r$ and global best $G$,
\begin{equation}
V_{r}^{(t+1)}=
\omega_{r,t}V_r^{(t)}
+\alpha_{r,t}\bigl(P_r^{(t)}-Q_r^{(t)}\bigr)
+\beta_{r,t}\bigl(G^{(t)}-Q_r^{(t)}\bigr),
\tag{39}
\end{equation}
\begin{equation}
Q_r^{(t+1)}=Q_r^{(t)}+V_r^{(t+1)}.
\tag{40}
\end{equation}
If a numerical coordinate leaves its observed range
$[L_j,U_j]$, it is reset to the dataset mean $\overline X_j$ rather than
clipped:
\begin{equation}
q_{rcj}^{(t+1)}\leftarrow\overline X_j
\quad\text{if }q_{rcj}^{(t+1)}\notin[L_j,U_j].
\tag{41}
\end{equation}

\subsection{Categorical update}

For a categorical coordinate, inequality is converted to a binary difference.
The velocity update is
\begin{equation}
V_{r}^{(t+1)}=
\omega_{r,t}V_r^{(t)}
+\alpha_{r,t}\I(P_r^{(t)}\ne Q_r^{(t)})
+\beta_{r,t}\I(G^{(t)}\ne Q_r^{(t)}).
\tag{42}
\end{equation}
The transfer probability is
\begin{equation}
T(v)=\left|\frac{v}{\sqrt{v^2+1}}\right|.
\tag{43}
\end{equation}
With $U\sim\operatorname{Uniform}(0,1)$, the coordinate is replaced by an
independent uniformly sampled category if $U<T(V_r^{(t+1)})$; otherwise it is
retained.  A newly sampled category may equal the current category.

\subsection{Feasible evaluation and best updates}

Every particle is assigned by (7).  It is discarded when any $n_c<k$.
For a feasible particle, the score is (24).  Personal and global bests are
updated only under strict score improvement:
\[
P_r^{(t+1)}=
\begin{cases}
Q_r^{(t)},&F_\lambda(Q_r^{(t)})<F_\lambda(P_r^{(t)}),\\
P_r^{(t)},&\text{otherwise},
\end{cases}
\tag{44}
\]
\[
G^{(t+1)}\in\argmin_{(r,h)\in\mathcal E_t}F_\lambda(Q_r^{(h)}),
\qquad
\mathcal E_t=\{(r,h):h\le t,\ Q_r^{(h)}\text{ is hard-$k$ feasible}\}.
\tag{45}
\]
For population $P$ and iteration count $T$, a normal completed scalar run
performs
\begin{equation}
N_{\mathrm{eval}}=PT
\tag{46}
\end{equation}
particle evaluations.  The recorded $P=T=100$ setting gives $10{,}000$.
The feasible-evaluation percentage is
\begin{equation}
100\frac{N_{\mathrm{feasible}}}
{N_{\mathrm{feasible}}+N_{\mathrm{infeasible}}}.
\tag{47}
\end{equation}
If the initial scalar population contains no hard-$k$ feasible particle, the
run returns the status
\[
\code{no\_feasible\_particle}.
\]
No penalty-derived pseudo-result is invented.

\section{Additive Pareto multi-objective extension}

\subsection{Objective vector and dominance}

The extension searches the two objectives directly:
\begin{equation}
\min\;\mathbf f(x)=\bigl(IL(x),-H_{\mathrm{sum}}(x)\bigr).
\tag{48}
\end{equation}
Candidate $x$ dominates $y$, written $x\prec y$, exactly when
\begin{equation}
IL(x)\le IL(y),\qquad H_{\mathrm{sum}}(x)\ge H_{\mathrm{sum}}(y),
\tag{49}
\end{equation}
and at least one inequality is strict.  The archive is the set of discovered
hard-$k$ feasible candidates that are not dominated by another retained
candidate, subject to a finite archive-size limit.

For a particle's personal archive point, a new candidate replaces the
incumbent if it dominates it.  If neither dominates the other, replacement is
selected by a Bernoulli$(1/2)$ draw.  If the incumbent dominates the new
candidate, the incumbent remains.

\subsection{Archive uniqueness and crowding distance}

An archive candidate is treated as a duplicate when it has the same assignment
vector and numerically equal $IL$ and entropy sum.  For objective coordinate
$m\in\{1,2\}$, sort archive values as
$f_{(1)m}\le\cdots\le f_{(A)m}$.  Boundary entries receive infinite crowding
distance.  Each interior entry accumulates
\begin{equation}
d_{(r)}\mathrel{+}=
\frac{f_{(r+1)m}-f_{(r-1)m}}
{f_{(A)m}-f_{(1)m}}
\quad\text{when }f_{(A)m}>f_{(1)m}.
\tag{50}
\end{equation}
When the archive is oversized, an entry having the smallest finite $d$ is
removed; boundary points are preserved whenever finite alternatives exist.

Leader-selection weights are
\begin{equation}
\widetilde d_r=
\begin{cases}
d_r+10^{-12},&d_r<\infty,\\
1+\max\{d_h:d_h<\infty\},&d_r=\infty,
\end{cases}
\qquad
\pi_r=\frac{\widetilde d_r}{\sum_h\widetilde d_h}.
\tag{51}
\end{equation}
If every distance is infinite, the replacement value in (51) is $1$.  The
global leader is drawn from the archive with probability $\pi_r$ and is then
used in the same PSO updates (39)--(43).

\subsection{Initialization attempts}

If no feasible archive member is found, a fresh independent population is
sampled, up to the declared initialization-attempt limit $A_0$ (20 in the
recorded Pareto studies).  Iteration count advances only after a nonempty
archive exists.  Exhausting all attempts produces
\code{no\_feasible\_particle}, which remains an algorithmic search outcome.

\section{Equivalence classes and disclosure-risk mathematics}

\subsection{Released equivalence classes}

Define equality of released records by their complete released QID vectors:
\begin{equation}
i\sim h\iff x'_{ij}=x'_{hj}\quad\text{for every QID }j.
\tag{52}
\end{equation}
The equivalence classes are $E_1,\ldots,E_G$.  Because different configured
clusters can in principle release the same QID representatives, $G$ is
reported separately from configured $C$.

\subsection{Sample prosecutor/linkage risk}

For record $i\in E_{g(i)}$, the sample prosecutor-risk proxy is
\begin{equation}
r_i=\frac{1}{\card{E_{g(i)}}}.
\tag{53}
\end{equation}
For threshold $\tau$, the implementation uses the strict boundary
\begin{equation}
L_\tau=\sum_{i=1}^{n}\I(r_i>\tau),
\qquad \tau\in\{0.05,0.075,0.10\}.
\tag{54}
\end{equation}
It also reports
\[
r_{\max}=\max_i r_i,\qquad
\overline r=\frac1n\sum_i r_i,
\qquad
r_{\mathrm{med}}=\median_i(r_i),
\tag{55}
\]
and the exposed fraction $L_\tau/n$.

\subsection{Sample marketer risk}

Since every class contributes
$\card{E_g}(1/\card{E_g})=1$ expected correct match under the implemented
sample model,
\begin{equation}
R_{\mathrm{marketer}}
=\frac1n\sum_{g=1}^{G}\card{E_g}\frac1{\card{E_g}}
=\frac Gn,
\tag{56}
\end{equation}
and the expected number of correct matches is $G$.  Journalist risk is not
estimated because population QID frequencies and the sampling fraction are not
available.

\subsection{Homogeneity exposure}

For SA $j$ in class $E_g$, let
\begin{equation}
d_{gj}=\card{\{S_{ij}:i\in E_g\}}.
\tag{57}
\end{equation}
The records exposed to homogeneity for SA $j$ are counted by
\begin{equation}
A_j=\sum_{g=1}^{G}\card{E_g}\I(d_{gj}=1).
\tag{58}
\end{equation}
The record-level union across all SAs is
\begin{equation}
A_{\cup}=\sum_{g=1}^{G}\card{E_g}
\I\!\left(\min_{1\le j\le p_S}d_{gj}=1\right),
\qquad
h_{\cup}=\frac{A_{\cup}}n.
\tag{59}
\end{equation}
Thus a record is counted once in $A_{\cup}$ even if its class is homogeneous
for several SAs.

\section{Diversity measures}

\subsection{Distinct \texorpdfstring{$\ell$}{l}-diversity}

For SA $j$ the achieved worst-class distinct diversity is
\begin{equation}
\ell^{\mathrm{distinct}}_j=\min_{1\le g\le G}d_{gj}.
\tag{60}
\end{equation}
The analysis also counts classes satisfying $d_{gj}\ge2$ and $d_{gj}\ge3$,
and records in classes below those levels:
\begin{equation}
N_{j,<\ell}=\sum_{g=1}^{G}\card{E_g}\I(d_{gj}<\ell),
\qquad \ell\in\{2,3\}.
\tag{61}
\end{equation}

\subsection{Entropy-effective diversity}

Let $p_{gjs}$ be the within-equivalence-class SA probability.  The analysis
uses natural-log entropy
\begin{equation}
H^{(e)}_{gj}=-\sum_s p_{gjs}\ln p_{gjs}
\tag{62}
\end{equation}
and its effective number of states
\begin{equation}
\ell^{\mathrm{eff}}_{gj}=\exp(H^{(e)}_{gj})
=2^{H^{(2)}_{gj}}.
\tag{63}
\end{equation}
The reported per-SA minimum is
\begin{equation}
\ell^{\mathrm{eff}}_j=\min_g\ell^{\mathrm{eff}}_{gj}.
\tag{64}
\end{equation}
These diversity quantities are post-anonymization measurements, not additional
optimizer constraints.

\subsection{Zero-minimum-entropy count}

The privacy screens use
\begin{equation}
Z_H=\sum_{j=1}^{p_S}\I\!\left(H_j^{\min}\approx0\right),
\tag{65}
\end{equation}
where the implementation uses numerical closeness to zero.

\section{Before--after distribution mathematics}

Let $X_j=(x_{1j},\ldots,x_{nj})$ and
$X'_j=(x'_{1j},\ldots,x'_{nj})$ be aligned before and after columns.

\subsection{Cell changes, distinctness, and missingness}

Missing-to-missing is treated as unchanged.  Define
\begin{equation}
c_{ij}=\I\bigl(x_{ij}\ne x'_{ij}\bigr),
\qquad
C_j=\sum_i c_{ij},
\qquad
f^{\mathrm{change}}_j=\frac{C_j}{n}.
\tag{66}
\end{equation}
The tables also report before/after numbers of distinct states and missing
cells.  For a scope containing $p$ attributes, its missing fraction is
\begin{equation}
f_{\mathrm{miss}}=\frac{\sum_{j=1}^{p}m_j}{np}.
\tag{67}
\end{equation}

\subsection{Numerical location, spread, and shape}

For the $m$ non-missing values $y_1,\ldots,y_m$ of a numeric attribute,
\begin{equation}
\overline y=\frac1m\sum_{i=1}^{m}y_i,
\qquad
s=\sqrt{\frac{1}{m-1}\sum_{i=1}^{m}(y_i-\overline y)^2}.
\tag{68}
\end{equation}
The median, first quartile $Q_{0.25}$, third quartile $Q_{0.75}$, and
interquartile range
\begin{equation}
IQR=Q_{0.75}-Q_{0.25}
\tag{69}
\end{equation}
are reported.  NumPy's default linear quantile interpolation is used: after
sorting, let $h=(m-1)q$, $r=\lfloor h\rfloor$, and $\gamma=h-r$; then
\begin{equation}
Q_q=(1-\gamma)y_{(r+1)}+\gamma y_{(r+2)},
\tag{70}
\end{equation}
with the obvious endpoint convention.  For $m>2$, pandas' adjusted Fisher--
Pearson sample skewness is
\begin{equation}
G_1=\frac{m}{(m-1)(m-2)}
\sum_{i=1}^{m}\left(\frac{y_i-\overline y}{s}\right)^3.
\tag{71}
\end{equation}
The invoked pandas routine returns $0$ for a numerically constant sample and an
undefined value when $m<3$; these are its executed edge-case branches around
(71).
Every change column is computed as
\begin{equation}
\Delta T=T_{\mathrm{after}}-T_{\mathrm{before}}.
\tag{72}
\end{equation}

\subsection{Numerical distribution distances}

For empirical CDFs $F$ and $G$, the one-dimensional Wasserstein distance is
\begin{equation}
W_1(F,G)=\int_{-\infty}^{\infty}|F(t)-G(t)|\,dt.
\tag{73}
\end{equation}
It is normalized by the before-data range
\begin{equation}
W_{\mathrm{norm}}=
\frac{W_1(F,G)}{\max(X_j)-\min(X_j)}.
\tag{74}
\end{equation}
If the range is zero, the value is $0$ only when raw $W_1=0$; otherwise it is
undefined.  The two-sample Kolmogorov--Smirnov statistic is
\begin{equation}
D_{KS}=\sup_t|F(t)-G(t)|.
\tag{75}
\end{equation}

\subsection{Categorical distribution distances}

After taking the union of before/after states and treating missingness as the
explicit state \code{<MISSING>}, let $p=(p_v)$ and $q=(q_v)$.  Total variation
distance is
\begin{equation}
D_{TV}(p,q)=\frac12\sum_v|p_v-q_v|.
\tag{76}
\end{equation}
Let $m=(p+q)/2$.  With base-2 logarithms,
\begin{equation}
\JSD_2(p,q)=\frac12\KL_2(p\|m)+\frac12\KL_2(q\|m),
\tag{77}
\end{equation}
\begin{equation}
\KL_2(p\|m)=\sum_{v:p_v>0}p_v\log_2\frac{p_v}{m_v}.
\tag{78}
\end{equation}
The code reports Jensen--Shannon \emph{distance}
\begin{equation}
D_{JS}(p,q)=\sqrt{\JSD_2(p,q)}.
\tag{79}
\end{equation}

\subsection{Attribute-level distance and scope similarity}

Let $\operatorname{mean}_{\mathrm{obs}}$ denote the pandas row mean over the
available applicable components (its default is to skip an undefined component).
The per-attribute composite is
\begin{equation}
d_j^{\mathrm{attr}}=
\begin{cases}
\operatorname{mean}_{\mathrm{obs}}
  \{W_{\mathrm{norm},j},D_{KS,j}\},
&j\text{ numeric},\\[2mm]
\operatorname{mean}_{\mathrm{obs}}
  \{D_{TV,j},D_{JS,j}\},
&j\text{ categorical}.
\end{cases}
\tag{80}
\end{equation}
When both applicable components are defined, as in the recorded analysis, each
row mean is their ordinary one-half sum.
For a scope $\mathcal A$ (all attributes, QIDs, or SAs),
\begin{equation}
\overline d_{\mathcal A}
=\frac1{\card{\mathcal A}}\sum_{j\in\mathcal A}d_j^{\mathrm{attr}},
\qquad
\operatorname{similarity}_{\mathcal A}=1-\overline d_{\mathcal A}.
\tag{81}
\end{equation}
The implementation also reports the mean and maximum of each component metric,
ignoring non-applicable missing metric cells.

\section{Cluster-level descriptive summaries}

For configured cluster $c$, the released proportion is
\begin{equation}
\pi_c=\frac{n_c}{n},
\qquad \sum_{c=1}^{C}\pi_c=1.
\tag{82}
\end{equation}
Across the $p_S$ cluster-specific SA entropies, the cluster table reports
\begin{equation}
H_c^{\min}=\min_j H_{cj},\qquad
\overline H_c=\frac1{p_S}\sum_{j=1}^{p_S}H_{cj},\qquad
H_c^{\max}=\max_j H_{cj},
\tag{83}
\end{equation}
and
\begin{equation}
Z_c=\sum_{j=1}^{p_S}\I(H_{cj}\approx0).
\tag{84}
\end{equation}

\section{Cross-seed summaries}

For a configuration with $R$ completed comparison seeds and metric values
$y_1,\ldots,y_R$, the repository reports
\begin{equation}
\overline y=\frac1R\sum_{r=1}^{R}y_r,
\qquad
s_y=\sqrt{\frac1{R-1}\sum_{r=1}^{R}(y_r-\overline y)^2},
\tag{85}
\end{equation}
together with $\min_r y_r$ and $\max_r y_r$.  Sample standard deviation is
undefined when $R<2$.  Completion percentage is
\begin{equation}
100\frac{R}{R_{\mathrm{planned}}}.
\tag{86}
\end{equation}
Worst-case screening uses maxima for loss and exposure and minima for cluster
size, for example
\begin{equation}
IL_{\max}=\max_r IL_r,
\quad
A_{\cup,\max}=\max_r A_{\cup,r},
\quad
n_{\min,\min}=\min_r\min_c n_{cr}.
\tag{87}
\end{equation}

\section{Privacy eligibility and project screening}

For one run or one Pareto archive candidate, define the project eligibility
indicator
\begin{align}
P={}&\I\!\left(\min_c n_c\ge k\right)
\prod_{\tau\in\{0.05,0.075,0.10\}}\I(L_\tau=0)\notag\\
&\times\I\!\left(\frac{A_{\cup}}n\le0.20\right)
\I(Z_H\le1)\I(IL\le0.60).
\tag{88}
\end{align}
These numerical limits are project screening rules, not claims of universal
release safety.

For a configuration with comparison-seed set $\mathcal S$, all-seed
eligibility is
\begin{equation}
P_{\mathrm{all}}=
\I(\text{every planned run is recorded and completed})
\prod_{s\in\mathcal S}P_s.
\tag{89}
\end{equation}
For the Pareto workflow, each seed first chooses from its eligible archive:
\begin{equation}
x_s^*=\argmin_{x\in\mathcal A_s:P(x)=1}
\bigl(IL(x),-H_{\mathrm{sum}}(x),\operatorname{ord}_s(x)\bigr)
\tag{90}
\end{equation}
in lexicographic order, where $\operatorname{ord}_s$ is the stable input-row
order (the stored archive candidate index in the executed study).  Configuration-level
mean IL and entropy are then
computed from the $x_s^*$ values, and $P_{\mathrm{all}}=1$ requires an eligible
$x_s^*$ for every seed.

\section{Pareto-front and configuration-selection mathematics}

\subsection{Descriptive cross-configuration Pareto frontier}

Among eligible scalar-grid configurations, configuration $a$ dominates $b$ in
the descriptive frontier if
\begin{equation}
\overline{IL}_a\le\overline{IL}_b,\qquad
\overline H_a\ge\overline H_b,
\tag{91}
\end{equation}
with at least one strict inequality.  This label is descriptive and does not
replace the declared automatic ranking.

\subsection{Default Pareto-study selection}

With zero information-loss margin, the Pareto study ranks all
$P_{\mathrm{all}}=1$ configurations lexicographically by
\begin{equation}
\left(\overline{IL},\ h_{\cup,\max},\ -k,\ C\right),
\tag{92}
\end{equation}
so the preferred order is lower mean IL, lower maximum homogeneity fraction,
higher $k$, then smaller $C$.

\subsection{Information-loss-margin selection}

Let
\begin{equation}
\overline{IL}_{\mathrm{best}}
=\min_{a:P_{\mathrm{all},a}=1}\overline{IL}_a.
\tag{93}
\end{equation}
For declared absolute margin $m\ge0$, the retained set is
\begin{equation}
\mathcal W_m=\left\{a:P_{\mathrm{all},a}=1,
\ \overline{IL}_a-\overline{IL}_{\mathrm{best}}
\le m+10^{-12}\right\}.
\tag{94}
\end{equation}
When $m>0$, the winner in $\mathcal W_m$ minimizes lexicographically
\begin{equation}
\left(h_{\cup,\max},\ \overline{IL},\ -\overline H,\ -k,\ C\right).
\tag{95}
\end{equation}
The confirmation protocol uses $m=0.01$.

\subsection{Scalar dense-range and joint-grid tie breakers}

The fixed-$k$, fixed-$\lambda$ dense-cluster workflow ranks eligible settings by
\begin{equation}
\left(\overline{IL},\ h_{\cup,\max},\ C\right).
\tag{96}
\end{equation}
The scalar joint $k/C/\lambda$ grid ranks them by
\begin{equation}
\left(\overline{IL},\ h_{\cup,\max},\ -k,\ C,\ -\lambda\right).
\tag{97}
\end{equation}
All selection is withheld until every declared comparison outcome is recorded.

\subsection{Exploratory seed-screen selection}

The historical seed screen for one fixed scalar configuration selected
\begin{equation}
s^*=\argmin_{s\in\mathcal S_{\mathrm{completed}}}
\bigl(F_\lambda(s),s\bigr),
\tag{98}
\end{equation}
that is, minimum scalar objective with smaller seed as the exact tie breaker.
This was an exploratory local-screen rule.

\section{Post-optimization \texorpdfstring{$\lambda$}{lambda} crossover analysis}

For archive candidate $a$, define its scalar line
\begin{equation}
F_a(\lambda)=IL_a-\lambda H_a,
\qquad \lambda\in[0,1].
\tag{99}
\end{equation}
For candidates $a$ and $b$ with $H_a\ne H_b$, their crossover is
\begin{equation}
\lambda_{ab}=\frac{IL_a-IL_b}{H_a-H_b}.
\tag{100}
\end{equation}
Only crossovers strictly inside $(0,1)$ are retained, rounded to 15 decimal
places, together with endpoints $0$ and $1$.  If consecutive boundaries are
$b_r<b_{r+1}$, their midpoint is
\begin{equation}
\lambda_r^{\mathrm{mid}}=\frac{b_r+b_{r+1}}2.
\tag{101}
\end{equation}
Within each seed $s$, the chosen archive member on that interval is
\begin{equation}
a_s(r)=\argmin_a
\bigl(F_{sa}(\lambda_r^{\mathrm{mid}}),\operatorname{index}(a)\bigr).
\tag{102}
\end{equation}
The interval's seed-mean values are
\begin{equation}
\overline{IL}_r=\frac1{\card{\mathcal S}}
\sum_{s\in\mathcal S}IL_{s,a_s(r)},
\qquad
\overline H_r=\frac1{\card{\mathcal S}}
\sum_{s\in\mathcal S}H_{s,a_s(r)},
\tag{103}
\end{equation}
and the reported midpoint objective is
\begin{equation}
\overline F_r^{\mathrm{mid}}
=\overline{IL}_r-\lambda_r^{\mathrm{mid}}\overline H_r.
\tag{104}
\end{equation}
Adjacent intervals having the same seed--candidate signature are merged.
After merging, (101) and (104) are recomputed using the merged endpoints; this
prevents a stale pre-merge midpoint value.

\section{Held-out candidate mathematics}

Comparison and held-out seeds must be disjoint:
\begin{equation}
s_{\mathrm{heldout}}\notin\mathcal S_{\mathrm{comparison}}.
\tag{105}
\end{equation}
The held-out Pareto run keeps only candidates satisfying (88).  If none does,
the run fails and another seed or candidate is not substituted.  Otherwise it
selects lexicographically
\begin{equation}
x_{\mathrm{heldout}}^*
=\argmin_{x:P(x)=1}\bigl(IL(x),-H_{\mathrm{sum}}(x),\operatorname{index}(x)\bigr).
\tag{106}
\end{equation}
This is evaluated once after aggregate selection.  It remains a technical
screening result rather than an unrestricted-release approval.

\section{Experiment-size arithmetic}

\subsection{General finite-grid count}

For a scalar grid with $k$ set $\mathcal K$, cluster interval
$[C_{\min},C_{\max}]$, lambda set $\Lambda$, and seed set $\mathcal S$, the
feasible endpoint for $k$ is
\begin{equation}
C_k^{\max}=\min\!\left(C_{\max},\left\lfloor\frac nk\right\rfloor\right).
\tag{107}
\end{equation}
The configuration and run counts are
\begin{equation}
N_{\mathrm{cfg}}^{\mathrm{scalar}}
=\card{\Lambda}\sum_{k\in\mathcal K}
\max\bigl(0,C_k^{\max}-C_{\min}+1\bigr),
\tag{108}
\end{equation}
\begin{equation}
N_{\mathrm{runs}}=N_{\mathrm{cfg}}\card{\mathcal S}.
\tag{109}
\end{equation}
The Pareto workflow removes $\lambda$ as an optimizer-grid dimension:
\begin{equation}
N_{\mathrm{cfg}}^{\mathrm{Pareto}}
=\sum_{k\in\mathcal K}
\max\bigl(0,C_k^{\max}-C_{\min}+1\bigr).
\tag{110}
\end{equation}
For explicitly listed pairs $\mathcal P\subset\mathcal K\times\mathcal C$,
\begin{equation}
N_{\mathrm{cfg}}=\card{\mathcal P},
\qquad
N_{\mathrm{runs}}=\card{\mathcal P}\card{\mathcal S}.
\tag{111}
\end{equation}

\subsection{Counts used in the recorded and configured work}

\begin{longtable}{@{}L{0.39\textwidth}L{0.55\textwidth}@{}}
\toprule
\textbf{Protocol} & \textbf{Arithmetic}\\
\midrule
\endfirsthead
\toprule
\textbf{Protocol} & \textbf{Arithmetic}\\
\midrule
\endhead
Pareto broad v1, $k\in\{10,15,20,25,30\}$, $C=1,\ldots,30$, ten seeds
& $5\times30=150$ configurations and $150\times10=1{,}500$ runs.\\
Pareto adaptive refinement, every integer $k=5,\ldots,30$, $C=2,\ldots,10$, ten seeds
& $26\times9=234$ configurations and $234\times10=2{,}340$ runs.\\
Pareto confirmation, three explicit $k/C$ pairs and ten fresh comparison seeds
& $3\times10=30$ runs.\\
Dense $k=10$, $C=1,\ldots,148$, ten seeds
& $148$ configurations and $148\times10=1{,}480$ runs; $148=\lfloor1485/10\rfloor$.\\
Scalar setup recorded in \code{settings/setup.tex}: every integer $k=5,\ldots,30$, $C=2,\ldots,20$, 20 lambda values, three seeds
& $26\times19\times20=9{,}880$ configurations and $9{,}880\times3=29{,}640$ runs.\\
Reusable scalar README example: eight listed $k$ values, $C=2,\ldots,20$, 20 lambda values, three seeds
& $8\times19\times20=3{,}040$ configurations and $3{,}040\times3=9{,}120$ runs.\\
Historical full sensitivity grid: four $k$ values, 11 cluster counts, five lambda values
& $4\times11\times5=220$ declared configurations.\\
\bottomrule
\end{longtable}

The held-out seed is excluded from all counts in the comparison-stage run
totals and contributes one separate final optimizer run only after selection.

\section{Release-readiness triage quantities}

The aggregate-only readiness module uses a few additional descriptive rules.
For non-missing text values in attribute $j$, it records maximum character
length, maximum token count, and mean token count
\begin{equation}
\overline t_j=\frac1{n_j^{\mathrm{obs}}}
\sum_{i:x_{ij}\text{ observed}}t_{ij}.
\tag{112}
\end{equation}
Potential free text is flagged when
\begin{equation}
\max_i\operatorname{characters}(x_{ij})\ge120
\quad\text{or}\quad
\max_i t_{ij}\ge8.
\tag{113}
\end{equation}
A phone-like value must contain between 7 and 15 digits inclusive and contain
only digits plus the permitted punctuation characters.  These are triage
heuristics, not disclosure guarantees.

\section{Important mathematical boundaries}

\begin{enumerate}
  \item Hard $k$ is enforced directly; the repository does not use the paper's
  ambiguous printed violation-penalty sign or invent a penalty multiplier.
  \item No unpublished normalization is applied to $IL$, entropy, or the
  scalar objective.
  \item Multiple-SA entropy is the sum of per-SA minimum cluster entropies in
  (P18)--(P19), not the minimum of within-cluster entropy sums.
  \item Linkage exposure uses the strict operator $r_i>\tau$.
  \item $\code{no\_feasible\_particle}$ describes a failed stochastic search,
  not proof that the mathematical feasible set is empty.
  \item PSO output is a best-found stochastic result under a recorded budget,
  not a proof of a global optimum.
  \item Passing (88) is an internal screen, not permission for unrestricted
  release.
  \item Objectives evaluated under different $\lambda$, schemas, SA counts,
  or preprocessing policies are not directly comparable.
\end{enumerate}

\appendix
\section{Implementation source map}

\begin{longtable}{@{}L{0.40\textwidth}L{0.54\textwidth}@{}}
\toprule
\textbf{Mathematics} & \textbf{Primary implementation source}\\
\midrule
\endfirsthead
\toprule
\textbf{Mathematics} & \textbf{Primary implementation source}\\
\midrule
\endhead
MO-OBAM deviations, information loss, entropy, hard $k$, objective, anonymization
& \code{mo\_obam/formulation.py}\\
Exact partitions, optimal fixed-membership centroids
& \code{mo\_obam/exact\_solver.py}\\
Profile weighting, compiled distance, scalar PSO updates
& \code{mo\_obam/pso\_solver.py}\\
Pareto dominance, archive, crowding, leader probabilities
& \code{mo\_obam/prti\_multiobjective\_solver.py}\\
Eligibility, seed aggregation, margin selection, lambda intervals
& \code{mo\_obam/prti\_multiobjective.py}\\
Held-out Pareto candidate rule
& \code{mo\_obam/prti\_multiobjective\_release.py}\\
Equivalence-class linkage and homogeneity measures
& \code{mo\_obam/privacy\_metrics.py}\\
Prosecutor/marketer risk, diversity, distribution distances and summaries
& \code{mo\_obam/prti\_analysis.py}\\
Preparation frequencies, cardinality, missingness, round-trip checks
& \code{mo\_obam/prti\_data.py}\\
Cross-seed means, sample deviations, feasible-evaluation percentage
& \code{mo\_obam/prti\_robustness.py}\\
Dense-range and joint scalar-grid selection
& \code{mo\_obam/prti\_dense\_clusters.py}\\
Scalar held-out screening
& \code{mo\_obam/prti\_release.py}\\
Cluster entropy summaries and transformation validation
& \code{mo\_obam/prti\_runner.py}\\
Text/identifier triage quantities
& \code{mo\_obam/prti\_release\_readiness.py}\\
Recorded mathematical narrative and implementation boundaries
& \code{PRTI\_MO\_OBAM\_DOCUMENTATION.md}\\
Delegated quantiles, sample skewness, Wasserstein/KS, and Jensen--Shannon distance
& NumPy \code{quantile}, pandas \code{Series.skew}, and SciPy
  \code{wasserstein\_distance}, \code{ks\_2samp}, and \code{jensenshannon}\\
\bottomrule
\end{longtable}

\section{Compact symbol table}

\begin{longtable}{@{}L{0.20\textwidth}L{0.73\textwidth}@{}}
\toprule
\textbf{Symbol} & \textbf{Meaning}\\
\midrule
\endfirsthead
\toprule
\textbf{Symbol} & \textbf{Meaning}\\
\midrule
\endhead
$n$ & number of records\\
$p_N,p_C,p_Q,p_S$ & numerical QIDs, categorical QIDs, all QIDs, and SAs\\
$C$ & configured number of clusters\\
$k$ & minimum records required in every configured cluster\\
$w_{ic}$ & binary membership of record $i$ in cluster $c$\\
$q_{cj}$ & QID representative for cluster $c$, attribute $j$\\
$d(i,c)$ & mixed squared-numeric plus categorical-mismatch distance\\
$IL$ & information-loss ratio\\
$H_{cj}$ & entropy of SA $j$ in cluster $c$\\
$H_j^{\min}$ & minimum cluster entropy for SA $j$\\
$H_{\mathrm{sum}}$ & sum of the $p_S$ per-SA minimum entropies\\
$\lambda$ & scalar entropy weight\\
$E_g$ & released QID equivalence class\\
$L_\tau$ & records with $1/\card{E_g}>\tau$\\
$A_j,A_\cup$ & per-SA and any-SA homogeneity-exposed record counts\\
$Z_H$ & SAs whose minimum cluster entropy is numerically zero\\
$P$ & one-run/candidate project eligibility indicator\\
$P_{\mathrm{all}}$ & all-comparison-seed eligibility indicator\\
\bottomrule
\end{longtable}

\end{document}
